\section{Discussion}
\textbf{What the noise tells us about Tamil.}
Vowel omission and key slips destroy the information the model needs.
Dropped vowels make the input genuinely ambiguous (a string like \emph{kankl} fits several Tamil words), so part of the remaining error is irreducible without context or a lexicon; top-5 accuracy and language-model reranking are therefore natural next steps rather than larger models.
By contrast, vowel-length and doubled-letter errors cost only 1 to 2 CER points, possibly because the Dakshina training data already contain many romanizations that collapse these distinctions~\cite{roark2020dakshina,wolfsonkin2019latin}; we did not test this.
It suggests that a Tamil-aware noise model should concentrate on vowel omission and keyboard errors.

\textbf{Data, not capacity.}
Increasing the BiLSTM from 1M to 8M parameters changes clean test CER by at most 0.25, while adding 60{,}000 pool words (A3 and CTRL) lowers it by 0.3 to 0.5 and the independent-set CER by about 1 to 1.5.
On this task the limiting factor is the amount and variety of romanized data.
This is consistent with the large gain that Aksharantar reports from mined data~\cite{madhani2023aksharantar}, and it motivates the larger-data fine-tuning described in \Cref{sec:final}.

\textbf{Back-transliteration: words or variants?}
Previous work on romanized language identification found that sampled spelling variants beat larger models~\cite{benton2025romanizedlid}.
Our control isolates the factors for transliteration: the clean-accuracy gain comes from extra \emph{words}, while spelling variants add a smaller robustness gain on mixed noise.
Because the variants come from a model trained on Dakshina, they reproduce Dakshina-style variation and not chat-style variation, so a larger effect may need variants that reflect real typing.

\textbf{Operator-specific robustness.}
Training on V-DROP gives a large gain on V-DROP and none on a held-out operator, as Belinkov and Bisk found for natural language noise~\cite{belinkov2018synthetic}.
Practical systems therefore need noise operators that cover the phenomena users produce, and the benchmark we release is a way to check coverage rather than a replacement for real data.

\textbf{Architectures.}
The conflict between Roark et al.\ and Kirov et al.\ is largely a matter of tuning and seeds: with 20 tuned trials per cell and five seeds the BiLSTM and the Transformer tie at 3M parameters, and the Transformer leads only under heavy noise.
Our tuning showed that Transformers are the more sensitive of the two to hyperparameters (\Cref{fig:tuning}), which supports the view that untuned comparisons can go either way~\cite{wu2021applying,moran2020effective}.
